\section{Tessuto emolinfopoietico}

\textbf{Connettivi liquidi}: tessuti in cui le cellule proprie del tessuto sono immerse in una matrice liquida.

% ---------------------------------------------------------------------------
\subsection{Il sangue}

Liquido vischioso, opaco, di colore rosso.
\begin{itemize}
  \item peso specifico: tra 1,055 e 1,065\,g/cm\textsuperscript{3};
  \item pH: tra 7,3 e 7,4;
  \item quantità totale: 7\% del peso corporeo, tra 5 e 6 litri.
\end{itemize}

È costituito da:
\begin{itemize}
  \item \textbf{matrice liquida} (sostanza intercellulare): il \textbf{plasma} (55\% del volume del sangue intero);
  \item \textbf{elementi figurati}: globuli rossi (99,9\%), globuli bianchi, piastrine.
\end{itemize}
Gli elementi figurati (45\% del volume) sono per la maggior parte prodotti dal \textbf{tessuto emopoietico}, che nell'adulto corrisponde al midollo osseo. Tra plasma ed eritrociti si interpone un sottile strato di leucociti e piastrine (meno dell'1\% del volume).

\rubrica{Funzioni}
\begin{itemize}
  \item \textbf{trasporto}: di $O_2$ e $CO_2$, elementi nutritivi (assorbiti a livello intestinale o prodotti nell'organismo), ormoni, anticorpi, enzimi, rifiuti metabolici;
  \item \textbf{regolazione}: della temperatura, dell'equilibrio acido-base, ionico e idrico;
  \item \textbf{difesa}: contro tossine e patogeni;
  \item \textbf{emostasi}: capacità di bloccare il flusso di sangue se viene lesionata una parte dei condotti in cui il sangue stesso scorre.
\end{itemize}

% ---------------------------------------------------------------------------
\subsection{Il plasma}

\begin{itemize}
  \item \textbf{acqua}: 92\%;
  \item \textbf{plasmaproteine} (7\%): albumina, globuline, fibrinogeno; prodotte dal fegato (90\%), dalle plasmacellule e da ghiandole endocrine;
  \item \textbf{sostanze inorganiche}: sodio, cloro, calcio, potassio, iodio, bicarbonato;
  \item \textbf{sostanze organiche}: materiali nutritivi e residui del metabolismo cellulare.
\end{itemize}

\rubrica{Funzioni delle proteine plasmatiche}
\begin{itemize}
  \item regolazione della pressione osmotica del sangue;
  \item regolazione del pH del sangue;
  \item intervento nei meccanismi di emostasi;
  \item intervento nell'immunità umorale;
  \item trasporto di varie molecole organiche ed inorganiche;
  \item riserva di materiale proteico.
\end{itemize}

\rubrica{Albumina}
\begin{itemize}
  \item 60\% delle proteine plasmatiche;
  \item sintetizzata dal fegato e da qui immessa nel sangue;
  \item mantiene la pressione osmotica del plasma;
  \item trasporta piccole molecole: calcio, bilirubina, acidi grassi liberi, amminoacidi essenziali, ormoni tiroidei (tiroxina), ormoni steroidei (cortisolo).
\end{itemize}

\rubrica{Globuline}
\begin{itemize}
  \item 35\% delle proteine plasmatiche;
  \item \textbf{lipoproteine} (HDL, LDL): trasportano lipidi e vitamine liposolubili;
  \item \textbf{protrombina}: fattore di coagulazione;
  \item \textbf{eritropoietina}: ormone che stimola la proliferazione degli eritrociti;
  \item \textbf{transferrina}: veicola il ferro;
  \item \textbf{immunoglobuline}: anticorpi;
  \item globuline di trasporto di ioni.
\end{itemize}

\rubrica{Fibrinogeno}
\begin{itemize}
  \item 5\% delle proteine plasmatiche;
  \item prodotto dal fegato;
  \item ruolo importante nel processo di emostasi.
\end{itemize}

% ---------------------------------------------------------------------------
\subsection{Elementi figurati}

\begin{table}[htbp]
  \centering
  \small
  \renewcommand{\arraystretch}{1.3}
  \begin{tabularx}{\textwidth}{|l|c|c|c|X|}
    \hline
    \textbf{Elemento} & \textbf{Diametro} & \textbf{Numero/mm\textsuperscript{3}} & \textbf{Vita media} & \textbf{Funzione} \\
    \hline
    Eritrociti & 7--8\,$\mu$m & 4--6 milioni & 100--120 giorni & trasporto di $O_2$ e $CO_2$ \\
    \hline
    Neutrofili & 12--14\,$\mu$m & 3000--7000 & 6 ore--pochi giorni & distruzione dei batteri per fagocitosi \\
    \hline
    Eosinofili & 12--15\,$\mu$m & 100--400 & $\sim$5 giorni & spengono le risposte allergiche, uccidono i parassiti \\
    \hline
    Basofili & 10--14\,$\mu$m & 20--50 & poche ore--pochi giorni & rilasciano istamina e altri mediatori dell'infiammazione \\
    \hline
    Linfociti & 5--17\,$\mu$m & 1500--3000 & da ore ad anni & risposta immunitaria (attacco diretto, cellule T; anticorpi, cellule B) \\
    \hline
    Monociti & 14--24\,$\mu$m & 100--700 & mesi & fagocitosi; nei tessuti si sviluppano in macrofagi \\
    \hline
    Piastrine & 2--4\,$\mu$m & 150.000--500.000 & 5--10 giorni & sigillano piccole lesioni dei vasi, coagulazione \\
    \hline
  \end{tabularx}
\end{table}

% ---------------------------------------------------------------------------
\subsection{Midollo osseo ed emopoiesi}

Il midollo osseo, alloggiato tra le trabecole ossee, contiene cellule progenitrici e cellule adipose; viene prelevato mediante aspirato osseo o biopsia osteomidollare.

Presenza di \textbf{emocitoblasti}, che differenziandosi danno origine ai progenitori degli elementi figurati del sangue:
\begin{itemize}
  \item \textbf{cellule staminali mieloidi}: differenziandosi danno origine ai globuli rossi e a diversi tipi di globuli bianchi;
  \item \textbf{cellule staminali linfoidi}: differenziandosi danno origine ai linfociti.
\end{itemize}

Dalla \textbf{cellula staminale pluripotente} del midollo osseo derivano due linee, ciascuna responsabile di distinti processi maturativi:
\begin{center}
\begin{forest} classalbero
[Staminale pluripotente
  [Staminale mieloide
    [Eritropoiesi]
    [Trombopoiesi]
    [Leucocitopoiesi]
  ]
  [Staminale linfoide
    [Linfociti B]
    [Linfociti T]
    [Cellule NK]
  ]
]
\end{forest}
\end{center}
\begin{itemize}
  \item \textbf{eritropoiesi}: proeritroblasto $\rightarrow$ eritroblasto $\rightarrow$ reticolocita $\rightarrow$ eritrocita;
  \item \textbf{trombopoiesi}: megacarioblasto $\rightarrow$ megacariocita $\rightarrow$ trombociti;
  \item \textbf{leucocitopoiesi}: mieloblasti $\rightarrow$ granulociti (neutrofili, eosinofili, basofili); monoblasti $\rightarrow$ monociti $\rightarrow$ macrofagi;
  \item la linea \textbf{linfoide} genera i linfociti B, T e NK.
\end{itemize}

% ---------------------------------------------------------------------------
\subsection{Globuli rossi (eritrociti o emazie)}

\begin{itemize}
  \item privi di nucleo;
  \item forma di dischi biconcavi, con grande rapporto superficie/volume;
  \item contenenti emoglobina;
  \item vita media di circa 120 giorni;
  \item numero da 4,2 a 6,5 milioni per mm\textsuperscript{3} di sangue;
  \item formano ``pile'' con facilitazione del loro flusso nei capillari, transitando anche attraverso capillari di 4\,$\mu$m di diametro;
  \item prodotti per \textbf{eritropoiesi} nel midollo osseo.
\end{itemize}

\rubrica{Eritropoiesi}
Catena differenziativa:
$$emocitoblasta \rightarrow \text{cellula staminale mieloide} \rightarrow proeritroblasto \rightarrow eritroblasto \rightarrow reticolocita \rightarrow eritrocita$$
\begin{itemize}
  \item \textbf{emocitoblasta}: progenitore totalmente indifferenziato;
  \item \textbf{cellula staminale mieloide}: cellula leggermente differenziata;
  \item \textbf{proeritroblasto}: progenitore dei globuli rossi;
  \item \textbf{eritroblasto}: deriva dal precedente;
  \item \textbf{reticolocita}: nucleo assente, ma organuli cellulari presenti;
  \item \textbf{eritrocita}: assenti anche gli organuli; è il globulo rosso maturo.
\end{itemize}
L'attività di eritropoiesi dipende da:
\begin{itemize}
  \item \textbf{ormoni}: eritropoietina, tiroxina, androgeni, ormone somatotropo;
  \item presenza di amminoacidi, ferro e vitamine (B6, B12 e acido folico).
\end{itemize}

\rubrica{Emoglobina}
\begin{itemize}
  \item 95\% delle proteine intracellulari del globulo rosso;
  \item contenuto di emoglobina del sangue intero: 14--18\,g/100\,ml nel maschio, 12--16\,g/100\,ml nella femmina;
  \item trasportatore dei gas;
  \item 4 subunità proteiche globulari; ciascuna porta un \textbf{gruppo eme}, che lega uno ione ferro ($Fe^{2+}$);
  \item quando lega ossigeno forma \textbf{ossiemoglobina};
  \item quando lega anidride carbonica forma \textbf{carbossiemoglobina};
  \item riciclo dell'emoglobina quando il globulo rosso muore.
\end{itemize}

\rubrica{Omeostasi di eritrociti ed emoglobina}
Gli eritrociti prodotti nel midollo osseo vengono riversati nel circolo sanguigno (veicolazione dei gas). Il midollo osseo può costruire cellule grazie alle risorse recuperate dal metabolismo dopo l'assorbimento e la digestione dei cibi. Nella \textbf{milza} avviene la demolizione dei globuli rossi, con successivo riutilizzo nel fegato degli amminoacidi e del ferro e con lo smaltimento dei prodotti di rifiuto.

% ---------------------------------------------------------------------------
\subsection{Gruppi sanguigni}

\rubrica{Sistema ABO}
Presenza o meno sul globulo rosso di antigeni di superficie (\textbf{agglutinogeni}) capaci di scatenare una risposta immunitaria. Gli anticorpi corrispondenti sono le \textbf{agglutinine}.

\begin{table}[htbp]
  \centering
  \small
  \renewcommand{\arraystretch}{1.3}
  \begin{tabular}{|c|c|c|c|c|}
    \hline
    \textbf{Antigene} & \textbf{Gruppo} & \textbf{Agglutinina} & \textbf{Riceve da} & \textbf{Dona a} \\
    \hline
    A        & A  & beta (anti-B)      & A e 0 & A e AB \\
    \hline
    B        & B  & alfa (anti-A)      & B e 0 & B e AB \\
    \hline
    entrambi & AB & nessuna            & tutti & AB \\
    \hline
    nessuno  & 0  & entrambe           & 0     & tutti \\
    \hline
  \end{tabular}
\end{table}

Il gruppo AB è \textit{ricevente universale}; il gruppo 0 è \textit{donatore universale}. Trasfusioni tra gruppi incompatibili provocano \textbf{emoagglutinazione}.

\rubrica{Fattore Rh}
Presenza, nei globuli rossi di una scimmia (\textit{Macacus rhesus}) e dell'uomo, di un nuovo antigene: il \textbf{fattore Rh}.
\begin{itemize}
  \item \textbf{Rh+}: il fattore è presente;
  \item \textbf{Rh$-$}: il fattore manca.
\end{itemize}
\textbf{Incompatibilità materno-fetale}: feto Rh+ con madre Rh$-$ e padre Rh+ $\rightarrow$ \textbf{malattia emolitica del neonato}. Alcuni globuli rossi Rh+ del feto (con l'antigene D) passano attraverso la placenta nel sangue materno; la madre produce anticorpi anti-D; in una gravidanza successiva gli anticorpi anti-D attraversano la placenta ed entrano nel sangue del feto, provocando l'emolisi del sangue Rh+.

% ---------------------------------------------------------------------------
\subsection{Piastrine (trombociti)}

Elementi morfologici anucleati, fondamentali per il \textbf{processo di coagulazione del sangue}.
\begin{itemize}
  \item si formano nel midollo osseo a partire da \textbf{megacariociti}, il cui citoplasma si suddivide in piccoli frammenti circondati da membrana (sino a 4000 piastrine per cellula);
  \item da 150.000 a 450.000/mm\textsuperscript{3};
  \item sopravvivono in circolo tra 8 e 12 giorni;
  \item formazione regolata da \textbf{trombopoietina}, interleuchina 6 e fattori di crescita specifici.
\end{itemize}

\rubrica{Struttura}
Membrana plasmatica, microtubuli, sistema tubulare denso, mitocondri, glicogeno, lisosomi (granuli lambda), granuli alfa, granuli delta, tubuli che si aprono in superficie.

% ---------------------------------------------------------------------------
\subsection{Emostasi}

\rubrica{Fase vascolare}
\begin{itemize}
  \item lesione del vaso sanguigno;
  \item contrazione delle fibre muscolari lisce del vaso;
  \item rilascio di sostanze chimiche ed ormoni locali.
\end{itemize}

\rubrica{Fase piastrinica}
\begin{itemize}
  \item le piastrine aderiscono alle superfici esposte della parete del vaso;
  \item modificazione della propria forma;
  \item creazione del \textbf{tappo piastrinico}.
\end{itemize}

\rubrica{Fase coagulativa}
\begin{itemize}
  \item fibrinogeno convertito in \textbf{fibrina};
  \item fattori rilasciati dalle piastrine e dalle cellule endoteliali interagiscono con i fattori della coagulazione;
  \item formazione del \textbf{coagulo sanguigno}.
\end{itemize}

\rubrica{Cascata coagulativa}
Due vie confluiscono nell'attivazione del fattore X:
\begin{itemize}
  \item \textbf{via intrinseca} (difetto cellulare tessutale): attiva il fattore III;
  \item \textbf{via estrinseca} (contatto con superficie non endoteliale vasale): fattore XII $\rightarrow$ fattore XI $\rightarrow$ fattore IX ($Ca^{2+}$).
\end{itemize}
$$\text{fattore X} \rightarrow \text{attivatore di protrombina} \rightarrow \text{trombina} \rightarrow \text{conversione fibrinogeno in fibrina}$$
L'attivatore di protrombina comprende fattori V--VI, complesso fattore VIII, fattore piastrine 3, $Ca^{2+}$. La trombina, insieme al fattore XIII attivato, converte il fibrinogeno in monomeri di fibrinogeno $\rightarrow$ \textbf{rete di fibrina}. In seguito la \textbf{plasmina} scioglie lentamente la fibrina $\rightarrow$ lisi del coagulo e riparazione della parete del vaso.

% ---------------------------------------------------------------------------
\subsection{Globuli bianchi (leucociti)}

Cellule rotondeggianti nucleate, localizzate nel tessuto connettivo propriamente detto o negli organi linfoidi.
\begin{itemize}
  \item tra 4.000 e 10.000 per mm\textsuperscript{3} di sangue;
  \item \textbf{leucopenia}: $<$4.000/mm\textsuperscript{3}; \textbf{leucocitosi}: $>$10.000/mm\textsuperscript{3};
  \item movimento ameboide, con migrazione al di fuori del circolo ematico;
  \item \textbf{chemotassi}: attivazione da parte di specifici stimoli chimici;
  \item fagocitosi;
  \item lisosomi o vescicole di secrezione presenti nei granulociti.
\end{itemize}

\rubrica{Formula leucocitaria}
Rapporti percentuali tra le classi: neutrofili 50--70\%; linfociti 20--40\%; monociti 3--8\%; eosinofili 2--4\%; basofili 0,5--1\%.

% ---------------------------------------------------------------------------
\subsection{Granulociti}

Leucociti polimorfonucleati: i nuclei presentano numerosi lobi. In base all'affinità per i coloranti si distinguono tre tipi.

\rubrica{Neutrofili}
50--70\% del totale. Difesa dell'organismo contro le infezioni, con funzione fagocitaria (prima linea di difesa).
\begin{itemize}
  \item i granuli sono lisosomi che contengono idrolasi e lisozima;
  \item sequenza d'azione: attivazione $\rightarrow$ chemotassi (migrazione direzionata) $\rightarrow$ fagocitosi $\rightarrow$ uccisione del patogeno ingerito $\rightarrow$ secrezione degli enzimi lisosomiali;
  \item nel connettivo degenerano e vengono degradati dai macrofagi.
\end{itemize}

\rubrica{Eosinofili}
2--4\% del totale. Aumentano in presenza di affezioni allergiche e di parassitosi.
\begin{itemize}
  \item quando stimolati rilasciano il contenuto dei granuli mediante esocitosi (bersagli troppo grossi per essere fagocitati);
  \item capacità chemotattiche;
  \item capacità fagocitiche verso immunocomplessi e microrganismi.
\end{itemize}

\rubrica{Basofili}
0,5--1\% del totale. Intervengono nel processo infiammatorio; scarsa attività fagocitica.
\begin{itemize}
  \item i granuli contengono \textbf{eparina} (anticoagulante) e \textbf{istamina} (vasodilatatore);
  \item coinvolti nello \textit{shock anafilattico}: un antigene introdotto nell'organismo $\rightarrow$ produzione di IgE; a una riesposizione allo stesso antigene, questo si lega alle IgE presenti sui basofili $\rightarrow$ degranulazione $\rightarrow$ rilascio di istamina (vasodilatatore), leucotrieni (azione broncospastica) e altri fattori $\rightarrow$ vasodilatazione con ipotensione, perdita di liquidi dai capillari, formazione di edemi, ipersecrezione mucosa, broncocostrizione.
\end{itemize}

% ---------------------------------------------------------------------------
\subsection{Leucociti non granulari}

\rubrica{Monociti}
2--8\% del totale. Insieme ai neutrofili costituiscono la prima linea di difesa dell'organismo.
\begin{itemize}
  \item cellule molto grandi, con nucleo ``a ferro di cavallo'';
  \item movimenti ameboidi;
  \item nello spazio tissutale si differenziano in \textbf{macrofagi};
  \item attività fagocitaria.
\end{itemize}

\rubrica{Linfociti}
20--30\% del totale. Prodotti nei linfonodi, nella milza e in altre strutture linfatiche; presenti nel circolo sanguigno e nei tessuti linfatici. Tre classi:
\begin{itemize}
  \item \textbf{linfociti T}: dal midollo osseo raggiungono il timo; immunità cellulo-mediata;
  \item \textbf{linfociti B}: immunità umorale (plasmacellule, anticorpi);
  \item \textbf{linfociti NK} (\textit{natural killer}): sorveglianza immunitaria.
\end{itemize}

% ---------------------------------------------------------------------------
\subsection{Tessuto linfatico e immunità}

Il sistema linfatico è suddiviso in 3 compartimenti:
\begin{enumerate}
  \item \textbf{comparto delle cellule staminali}, nel midollo osseo;
  \item \textbf{organi linfatici primari} (centrali), in cui i linfociti assumono la competenza: timo per i linfociti T e midollo osseo per i linfociti B (negli uccelli, borsa di Fabrizio);
  \item \textbf{organi linfatici secondari} (periferici), in cui avviene l'incontro antigene-anticorpo (Ag-Ab): linfonodi, milza, sistema linfatico delle mucose.
\end{enumerate}
Il fluido interstiziale in eccesso ritorna al sangue attraverso i \textbf{vasi linfatici} (dotto linfatico destro e dotto toracico).

\rubrica{Linfonodo}
Organizzato in capsula, trabecole, seni (sottocapsulare, midollare), zona corticale superficiale (linfociti B), zona corticale profonda (linfociti T), cordoni midollari (linfociti B e plasmacellule); vaso afferente ed efferente, arteria e vena all'ilo.

\rubrica{Milza}
Presenta capsula e trabecole; \textbf{polpa rossa} e \textbf{polpa bianca}, seno venoso, arteria e vena splenica all'ilo.

\rubrica{Immunità innata}
\begin{enumerate}
  \item barriere di protezione del corpo;
  \item sistema dei fagociti $\rightarrow$ macrofagi (neutrofili ed eosinofili);
  \item cellule ad attività citotossica naturale (linfociti NK) contro cellule tumorali e cellule infettate da virus.
\end{enumerate}

\rubrica{Immunità acquisita}
Caratterizzata da \textbf{specificità} e \textbf{memoria immunologica}.
\begin{itemize}
  \item \textbf{immunità umorale}: mediata dai linfociti B e dai loro prodotti di secrezione, gli anticorpi; difesa nei confronti dei microbi extracellulari e dei loro prodotti (le tossine);
  \item \textbf{immunità cellulo-mediata}: mediata dai linfociti T e dai loro prodotti di secrezione, le citochine; difesa contro microbi intracellulari (virus e batteri che sopravvivono e proliferano all'interno dei fagociti e di altre cellule, inaccessibili agli anticorpi circolanti).
\end{itemize}

\rubrica{Anticorpi}
Gli anticorpi formano aggregati con molecole estranee, virus e batteri: gli aggregati di anticorpi ed antigeni sono fagocitati dalle cellule fagocitarie; speciali proteine nel sangue uccidono i batteri o i virus ricoperti di anticorpi.
